\documentclass[12pt]{article}

\usepackage[a4paper,margin=1.8cm]{geometry}
\usepackage{booktabs}
\usepackage{hyperref}

\begin{document}
	\title{\LARGE{INDIAN KNOWLEDGE SYSTEM PYQ}}
	\date{\nodate}
	\maketitle
	\chapter{QUICK STATS}
	\newpage
	
	\chapter{\LARGE{PYQS}}

	\begin{enumerate}
		\section{1-marks questions}
			\item What was the place value of the Babylonians' number system.
			\hfill{[CAT 25]}
			
			\item Name the person who give the first symbols for zero.What was that symbol? \hfill{[END 24]}
			
			\item What is the meaning of PARAM in PARAM-8000? \hfill{[END 22]}
			
			\item Write the equation of the author Chandrashshastra. \hfill{[END 22]}
			
			\item List top three Indian It companies in the decreasing order of their market capital. \hfill{[END 22]}
			
			\item List any three R\&D entities in India that has notable contribution in research in the domain of Computer Science. \hfill{[END 22]}
			
			
			
			
			
			
		\section{2-marks questions}
			\item Why do quantities more than four look vagure?Why does a person need to use counting for such quantites? \hfill{[CAT 25]}
			\item Which dynasty pioneered the preparation of wootz steel in India? \hfill[END 25]
			\item As per the current knowlwdge whoose text mentions fibonacci series? \hfill{[END 25]}
			\item Which are the top three Indian IT companies in term of market capital?\hfill{[END 25]}
			
			\item Who is the head of the committee that recommended estabblishing IITs in India and when this committee was formed? \hfill{[END 25]}
			
			\item List three objectives of establishing AICTE. \hfill{[END 25]}
			
			\item \textit{"We owe a lot to the Ancient Indians,teaching us how to count.Without which most modern scientific discoveries would have been impossible"}.Whoose quote is this? \hfill{[END 24]}
			
			\item What are the two uses of zero in number system?List them. \hfill{[END 24]}
			
			\item List any two objectives of the Indian Super Computing program. \hfill{[END 24]}
			
			\item When Indian IT industry was born? \hfill{[END 24]}
			
			\item Name the commitee and its chairman whoose recommendation guided the Govt. of India to extablish IITs. \hfill{[END 24]}
			
			\item Discuss any two methods ancient people used for counting cattles. \hfill{[CAT 24]}
			
			
			
			
			
		\section{3-marks questions}
			\item The Bakshali manuscript is famous for what?When and where was it found? \hfill{[CAT 25]}
			
			\item What do you understand about the Bakshali manuscript? \hfill{CAT 24}
		\section{4-marks questions}
			\item Explain why the Golden ratio is called the divine ratio. \hfill{[CAT 25]}
			\item Write about the contribution of Aryabhatta to Ancient Indian Knowlwdge. \hfill{[CAT 25]}
			
			\item Why was the development of counting methods necassary for humarn civilizations? \hfill{[CAT 24]}
			
			\item What is the golden ratio?Explain any three places/objects where you can observe the Golden Ratio. \hfill{[CAT 24]}
			
			\item What is Fibonacci Series?What is the Indian connection of this series? \hfill{[CAT 22]}
			
			\item List the names of three famous Indians and their contributin to Computer Science. \hfill{[CAT 22]}
			
			\item What are the key resons behinf the tremendous success of India in IT? \hfill{[CAT 22]}
			
			
			
			
		\section{4-marks questions}
			\item Discuss the concept of \textit{Golden Ration} with some examples \hfill{[CAT 22]}
			
			\item Write a short note on any two of the following:
			\begin{enumerate}
				\item CDAC
				\item Invention of zero
				\item Aryabhatta
			\end{enumerate}
			
			
			
			
			
			
			
			
			
			
		\section{5-marks questions}
		
		
		
		
		
		
		
		
		
		
		
		\section{10-marks questions}
		
			\item Describe the contribution of Aryabhatta to Indian Mathematics and astronomy. \hfill{[END 25]}
			
			\item Disscuss all the evidence of zero found in India. \hfill{[END 25]}
			
			\item Write a breif note on the following two ancient Indians who significantly contributed to Computing. \hfill{[END 25]}
			
			\begin{enumerate}
				\item Pingala and his contributution
				\item Brahmagupta and his contribution to mathematics.
			\end{enumerate}
			
			\item Explain R\&D ecosystem of Computer Science and Information technology in India. \hfill{[END 25]}
			
			\item Discuss about the ecosystem of funding for the development of science and technology in India. \hfill{[END 25]}
			
			\item What do you understand by \textit{sensing numbers} by humans and animals?Elaborate. \hfill{[END 24]}
			
			\item Put some light on the history and evidence of the number \textit{zero} in connection with Bharat. \hfill{[END 24]}
			
			\item Write a brief note on the following Indians who significantly contributed to Computing. \hfil{[END 24]}
			
			\begin{enumerate}
				\item Ajay Bhatt
				\item Abhay Bhushan
			\end{enumerate}
			
			\item Explain why AICTE was established. \hfill{[END 24]}
			
			\item What is the role of CDAC in creating supercomputing infrastructure in India?List the names of any three supercomputers developed in India. \hfill{[END 24]}
			
			\item What is the importance of zero in a number system? Elaborate on the invention of zero. \hfill{[END 22]}
			
			\item Write about 5 Indians and their contribution in Computing. \hfill{END 22}
			
			\item Elabrote the Super Computing program of India. \hfill{[END 22]}
			
			\item Write about three main Funding agency who fund research in computing. \hfill{[END 22]}
			
			\item Elabrote on the condition of Private sector in the field of computing in India. \hfill{[END 22]}
			
			
			
			
			
			
			
			
			
			
			
			
		\section{20-marks questions}
			\item Answer the following each carrying 10 points: \hfill{[END 25]}
			\begin{enumerate}
				\item Explain why the evolution of counting was neccessary to the evolution of human civilizations.
				
				\item Explain the presence of fibonacci numbers in the family tree of Honeybees.
			\end{enumerate}
			
			\item Answer the following each carrying 10 points: \hfill{[END 25]}
			
			\begin{enumerate}
				\item Discuss how India ecome global leader in Information Technology.
				\item Elabrote about the academic ecosystem of Computer Science and Information Technology in India.
			\end{enumerate}
			
			\item Answer the following each carrying 10 points: \hfill{[END 24]}
			
			\begin{enumerate}
				\item Explain the \textit{evolution of numbers} across different human civilizations.
				\item Explain the \textit{Golden Ratio} and discuss any two expressions of this ratio that are available in nature.
			\end{enumerate}
			
			\item Answer the following each carrying 10 points: \hfill{[END 24]}
			
			\begin{enumerate}
				\item Discuss briefly the top 5 IT companies of India.
				\item Explain how IITs were establieshed by Govt. of India.Also, elabrote on how they become leading technical institutes of the country.
			\end{enumerate}
			
			\item Elabrote on the ecosystem of computer science research and application development in India and it's mportant achivements. \hfill{[END 22]}
			
			\item What is the Indian connectio of Fibonacci Series?Elabrote on its presence in nature and its connection with Goldern Ratio. \hfill{[END 22]}
	\end{enumerate}
\end{document}